\chapter{\'Etale morphisms}
\label{I}

% original p. 1
To simplify the exposition, all the preschemes considered are assumed
to be locally noetherian, at least after no.~\ref{I.2}.

\section{Notions of differential calculus}
\label{I.1}

Let $X$ be a prescheme over~$Y$, and let $\Delta_{X/Y}$ or $\Delta$
be the diagonal morphism $X\to X\times_Y X$. It is an immersion,
hence a closed immersion of~$X$ into an open subset $V$ of~$X\times_Y X$.
Let $\cal{I}_X$ be the ideal of the closed subprescheme corresponding
to the diagonal in~$V$ (N.B. if one wants to proceed intrinsically,
without assuming $X$ separated over~$Y$---an assumption that would be
quite a joke---one should consider the set-theoretic inverse image
of~$\cal{O}_{X\times X}$ on~$X$, and denote by $\cal{I}_X$ the
augmentation ideal in the latter\ldots). The sheaf
$\cal{I}_X/\cal{I}_X^2$ can be regarded as a quasi-coherent sheaf
on~$X$; it is denoted by~$\mathit{\Omega}^1_{X/Y}$.
It is of finite type if $X\to Y$ is of finite type. It behaves well
with respect to extension of the base $Y'\to Y$. One also introduces
the sheaves $\cal{O}_{X\times_Y X}/\cal{I}_X^{n+1}=\cal{P}^n_{X/Y}$;
these are sheaves of \emph{rings} on~$X$, making $X$ into a prescheme
which one can denote by $\Delta_{X/Y}^n$ and call the
\emph{$n$th infinitesimal neighborhood of}~$X/Y$.
The corresponding sorites is utterly trivial, although rather
long\footnote{\cf EGA IV~16.3.}; it would be prudent to discuss it
only when one has something useful to say about it, in connection
with smooth morphisms.

\section{Quasi-finite morphisms}
\label{I.2}

\begin{proposition}
\label{I.2.1}
Let $A\to B$ be a local homomorphism (N.B. the rings are now
noetherian), and let $\goth{m}$ be the maximal ideal of~$A$.
The following conditions are equivalent:
\begin{enumerate}
\item[(i)] $B/\goth{m}B$ is finite-dimensional over~$k=A/\goth{m}$.
\item[(ii)] $\goth{m}B$ is an ideal of definition, and
$B/\goth{r}(B)=\kres(B)$ is an extension of~$k=\kres(A)$.
\item[(iii)] The completion $\widehat{B}$ is finite over the
completion $\widehat{A}$ of~$A$.
\end{enumerate}
\end{proposition}

One then says
% original p. 2
that $B$ \emph{is quasi-finite} over~$A$. A morphism $f\colon X\to Y$
is said to be quasi-finite at~$x$ (or the $Y$-prescheme $f$ is said
to be quasi-finite at~$x$) if $\cal{O}_x$ is quasi-finite
over~$\cal{O}_{f(x)}$. This also amounts to saying that $x$ is
\emph{isolated in its fiber}~$f^{-1}(x)$. A morphism is said to be
quasi-finite if it is so at every point\footnote{In EGA~II~6.2.3,
$f$ is also assumed to be of finite type.}.

\begin{corollary}
\label{I.2.2}
If $A$ is complete, quasi-finite is equivalent to finite.
\end{corollary}

One could give the usual sorites \textup{(i), (ii), (iii), (iv), (v),}
for quasi-finite morphisms, but this does not seem indispensable here.

\section{Unramified or net morphisms}
\label{I.3}

\begin{proposition}
\label{I.3.1}
Let $f\colon X\to Y$ be a morphism of finite type, $x\in X$,
$y=f(x)$. The following conditions are equivalent:
\begin{enumerate}
\item[(i)] $\cal{O}_x/\goth{m}_y\cal{O}_x$ is a finite separable
extension of~$\kres(y)$.
\item[(ii)] $\mathit{\Omega}^1_{X/Y}$ vanishes at~$x$.
\item[(iii)] The diagonal morphism $\Delta_{X/Y}$ is an open
immersion in a neighborhood of~$x$.
\end{enumerate}
\end{proposition}

For the implication (i)$\To$(ii), Nakayama immediately reduces
one to the case where $Y=\Spec(k)$, $X=\Spec(k')$, where this is
well known and, moreover, trivial from the definition of separability;
(ii)$\To$(iii) follows from a pleasant and easy characterization
of open immersions, using Krull; (iii)$\To$(i) because one is again
reduced to the case where $Y=\Spec(k)$ and the diagonal morphism
is an open immersion everywhere. One must then prove that $X$ is
finite with ring separable over~$k$; for this one is reduced to
the case where $k$ is algebraically closed. But then every closed
point of~$X$ is isolated (being identical to the inverse image of
the diagonal under the morphism $X\to X\times_k X$ defined by~$x$),
whence $X$ is finite. One may then suppose that $X$ consists of
a single point, with ring~$A$, so $A\otimes_k A\to A$ is an
isomorphism, whence $A=k$, qed.

\begin{definition}
\label{I.3.2}
\begin{enumerate}
\item[a)] One then says that $f$ is \emph{net}, or
\emph{unramified}, at~$x$, or that $X$ is net, or unramified,
at~$x$ over~$Y$.
\item[b)] Let $A\to B$ be a local homomorphism. One says that it
is \emph{net}, or \emph{unramified}, or that $B$ is a \emph{net},
or \emph{unramified}, local algebra over~$A$, if $B/\goth{r}(A)B$
is a finite separable extension of~$A/\goth{r}(A)$, \ie if
$\goth{r}(A)B=\goth{r}(B)$ and $\kres(B)$ is a separable extension
of~$\kres(A)$\footnote{\Cf the second thoughts in
III~\SourceRef{III.1.2}{1.2}.}.
\end{enumerate}
\end{definition}

\begin{remarks}
The fact
% original p. 3
that $B$ is net over~$A$ can already be recognized from the
completions of~$A$ and~$B$. Net implies quasi-finite.
\end{remarks}

\begin{corollary}
\label{I.3.3}
The set of points where $f$ is net is open.
\end{corollary}

\begin{corollary}
\label{I.3.4}
Let $X'$, $X$ be two preschemes of finite type over~$Y$, and
$g\colon X'\to X$ a $Y$-morphism. If $X$ is net over~$Y$, the
graph morphism $\Gamma_g\colon X'\to X\times_Y X$ is an open
immersion.
\end{corollary}

Indeed, it is the inverse image of the diagonal morphism
$X\to X\times_Y X$ under
\[
g\times_Y\id_{X'}\colon X'\times_Y X\to X\times_Y X.
\]

One can also introduce the annihilator ideal $\goth{d}_{X/Y}$
of~$\mathit{\Omega}^1_{X/Y}$, called the different ideal
of~$X/Y$; it defines a closed subprescheme of~$X$ whose underlying
set is the set of points where $X/Y$ is ramified, \ie not net.

\begin{proposition}
\label{I.3.5}
\begin{enumerate}
\item[(i)] An immersion is net.
\item[(ii)] The composite of two net morphisms is net.
\item[(iii)] Extension of the base in a net morphism gives another
net morphism.
\end{enumerate}
\end{proposition}

This can be seen equally well from (ii) or (iii) (the latter seems
more amusing to me). One can of course also be more precise by
giving pointwise statements; this is more general only in appearance
(except in the setting of definition~b)), and tedious. As usual,
one obtains corollaries:

\begin{corollaries}
\label{I.3.6}
\begin{enumerate}
\item[(iv)] A cartesian product of two net morphisms is another
net morphism.
\item[(v)] If $gf$ is net, $f$ is net.
\item[(vi)] If $f$ is net, $f_{\text{\textup{red}}}$ is net.
\end{enumerate}
\end{corollaries}

\begin{proposition}
\label{I.3.7}
Let $A\to B$ be a local homomorphism; suppose that the residue
field extension $\kres(B)/\kres(A)$ is trivial or that $\kres(A)$
is algebraically closed. For $B/A$ to be net, it is necessary and
sufficient that $\widehat{B}$ be (as an $\widehat{A}$-algebra)
a quotient of~$\widehat{A}$.
\end{proposition}

\begin{remarks}
\begin{enumerate}
\item[--] When the residue field extension is not assumed trivial,
one can reduce to this case by making a suitable finite flat
extension of~$A$ which kills the said extension.
\item[--] Give the example where $A$ is the local ring of an
ordinary double point of a curve, and $B$ that of a point of its
normalization: then $A\subset B$, $B$ is net over
% original p. 4
$A$ with trivial residue field extension, and
$\widehat{A}\to\widehat{B}$ is surjective but \emph{not injective}.
We shall therefore strengthen the notion of being net.
\end{enumerate}
\end{remarks}
